% !TeX root = ../概率论与数理统计笔记.tex
% 由主文档通过 \input 载入；请编译主文档。
\setcounter{chapter}{1} % 保留原书章号。
\chapter{查阅表、方法索引与综合学习检查}
\label{app:ps-reference}

本附录的表格与索引将在各章正文完成后同步整理，所有条目回指推导所在章节，并保留假设，避免形成脱离条件的公式清单。

\section{常见分布查阅表}
按分布名称、参数及范围、支撑、质量函数或密度、分布函数、均值、方差和典型生成机制组织。
统一几何及负二项的计数约定、伽马的率参数约定、正态的方差参数约定；存在多种参数化时列出转换关系。

\section{估计、区间与检验方法索引}
按推断目标、抽样机制、模型条件、所用统计量、精确或渐近分布以及正文位置组织。
将独立样本、配对样本、有限总体抽样分开列示，并标注边界参数和小样本情形需要重新检查的方法。

\section{综合例题与分层习题安排}
规划三条贯穿全书的综合问题：伯努利成功概率从分布到推断，正态测量误差从均值方差到回归，计数与等待从泊松分布到泊松过程。
各章习题依次设置概念辨析、基本计算、证明反例和综合建模；选定习题提供完整解答，其余提供提示或结果并明确标识。

\section{完成一项概率统计分析的检查顺序}
先确认研究问题、随机对象与数据获得方式，再写出模型及假设、确定目标量和方法、核对适用条件、执行推导或计算，最后解释结果与不确定性。
正文中的每个综合案例均按这一顺序组织，使分散的概念能够汇成可独立完成的分析过程。




